\section{Exercise 3}

Tệp \texttt{raw\_input.txt} chứa một dòng dữ liệu gồm năm trường thông tin ngăn cách bởi dấu phẩy. Trong vùng nhớ \texttt{.data}, chương trình khai báo sẵn năm nhãn kèm một mảng địa chỉ \texttt{label\_tab} lưu con trỏ tới từng nhãn theo đúng thứ tự. Nhờ mảng này, chương trình chỉ cần dựa vào số thứ tự của trường hiện tại để lấy thẳng nhãn tương ứng, thay vì phải viết nhiều lệnh rẽ nhánh để kiểm tra từng trường hợp:

\begin{codeblock}[language=MIPS]
.data
  file_in:    .asciiz "raw_input.txt"
  file_out:   .asciiz "formatted_result.txt"
  header:     .asciiz "-----Student personal information-----\n"
  msg_open:   .asciiz "Error opening file.\n"
  label_name: .asciiz "Name: "
  label_id:   .asciiz "ID: "
  label_addr: .asciiz "Address: "
  label_age:  .asciiz "Age: "
  label_rel:  .asciiz "Religion: "
  label_tab:  .word label_name, label_id, label_addr, label_age, label_rel
\end{codeblock}

Chương trình mở tệp đầu vào ở chế độ đọc, cấp phát vùng đệm $1000$ byte trên heap để chứa dữ liệu và lưu lại số byte thực tế đọc được từ \texttt{syscall 14}:

\begin{codeblock}[language=MIPS]
  li $v0, 13           # syscall 13: mở file
  la $a0, file_in      # đường dẫn file cần đọc
  li $a1, 0            # cờ 0: mở để đọc
  li $a2, 0            # mode bị bỏ qua
  syscall
  bltz $v0, open_error # báo lỗi nếu không mở được file đầu vào
  move $s0, $v0        # $s0 = file descriptor

  li $v0, 9            # syscall 9: cấp phát heap
  li $a0, 1000         # xin 1000 byte
  syscall
  move $s1, $v0        # $s1 = vùng đệm chứa nội dung file

  li $v0, 14           # syscall 14: đọc file
  move $a0, $s0        # đọc từ file vừa mở
  move $a1, $s1        # ghi vào vùng đệm heap
  li $a2, 1000         # đọc tối đa 1000 byte
  syscall
  move $s2, $v0        # $s2 = số byte thực đọc
\end{codeblock}

Sau khi đóng tệp nguồn, chương trình cấp phát một vùng heap khác rộng $2000$ byte để chứa chuỗi kết quả, trong đó \texttt{\$s3} giữ cố định địa chỉ đầu chuỗi và \texttt{\$s4} luôn trỏ tới vị trí cuối chuỗi hiện tại để nối tiếp dữ liệu. Trước khi vào vòng lặp, dòng tiêu đề và nhãn \texttt{"Name: "} được nối sẵn vào vùng đệm này:

\begin{codeblock}[language=MIPS]
  li $v0, 16         # syscall 16: đóng file
  move $a0, $s0      # descriptor cần đóng
  syscall

  li $v0, 9          # syscall 9: cấp phát heap cho chuỗi kết quả
  li $a0, 2000       # xin 2000 byte
  syscall
  move $s3, $v0      # $s3 = địa chỉ đầu chuỗi kết quả
  move $s4, $v0      # $s4 = con trỏ ghi, chạy theo từng ký tự

  la $a0, header     # chuỗi cần nối phải nằm trong $a0
  jal append_str     # nối dòng tiêu đề
  la $a0, label_name # chuỗi nhãn của trường đầu tiên
  jal append_str     # nối "Name: "

  move $t1, $s1      # $t1 = con trỏ đọc trong vùng đệm
  add $t3, $s1, $s2  # $t3 = địa chỉ kết thúc dữ liệu đọc được
  li $t0, 0          # $t0 = chỉ số trường hiện tại
\end{codeblock}

Vòng lặp \texttt{copy\_loop} duyệt từng byte trong vùng đệm: nếu gặp dấu phẩy thì chuyển sang trường kế tiếp; nếu gặp ký tự xuống dòng (\texttt{'\textbackslash n'}), ký tự kết thúc chuỗi (\texttt{'\textbackslash 0'}) hoặc hết dữ liệu thì dừng lại:

\begin{codeblock}[language=MIPS]
copy_loop:                 # chép từng ký tự cho tới hết dòng
  beq $t1, $t3, copy_done  # đọc hết dữ liệu thì dừng

  lb $t2, 0($t1)           # $t2 = ký tự hiện tại

  beq $t2, '\n', copy_done # gặp xuống dòng thì dừng
  beqz $t2, copy_done      # gặp ký tự kết thúc chuỗi thì dừng
  beq $t2, ',', next_field # gặp dấu phẩy: sang trường kế tiếp

  move $a0, $t2            # tham số của append_char là ký tự
  jal append_char          # chép ký tự sang chuỗi kết quả

  addi $t1, $t1, 1         # tiến con trỏ đọc
  j copy_loop
\end{codeblock}

Tại nhãn \texttt{next\_field}, chương trình ghi một ký tự xuống dòng thay cho dấu phẩy, nhân chỉ số trường với $4$ (bằng lệnh dịch trái \texttt{sll} $2$ bit) để tra địa chỉ con trỏ nhãn trong \texttt{label\_tab}, nạp địa chỉ nhãn qua \texttt{lw} và nối vào chuỗi đích:

\begin{codeblock}[language=MIPS]
next_field:         # hết một trường: xuống dòng và in nhãn kế tiếp
  addi $t0, $t0, 1  # sang trường kế tiếp
  li $a0, '\n'      # ký tự xuống dòng
  jal append_char   # xuống dòng trước nhãn mới
  la $t4, label_tab # $t4 = địa chỉ bảng nhãn
  sll $t5, $t0, 2   # $t5 = 4 * chỉ số trường
  add $t5, $t4, $t5 # $t5 = &label_tab[i]
  lw $a0, 0($t5)    # $a0 = nhãn của trường thứ i
  jal append_str    # nối nhãn vào chuỗi kết quả
  addi $t1, $t1, 1  # bỏ qua dấu phẩy
  j copy_loop
\end{codeblock}

Khi kết thúc vòng lặp, chuỗi kết quả được bổ sung ký tự xuống dòng và ký tự kết thúc chuỗi, sau đó in ra console:

\begin{codeblock}[language=MIPS]
copy_done:         # đã chép xong cả dòng
  li $a0, 10       # mã ASCII của ký tự xuống dòng
  jal append_char  # kết thúc dòng cuối
  sb $zero, 0($s4) # đặt ký tự kết thúc chuỗi

  li $v0, 4        # syscall 4: in chuỗi
  move $a0, $s3    # in toàn bộ chuỗi kết quả
  syscall
\end{codeblock}

Tiếp theo, chương trình mở tệp \texttt{formatted\_result.txt} ở chế độ ghi và ghi toàn bộ nội dung vừa định dạng xuống tệp qua \texttt{syscall 15}. Độ dài dữ liệu cần ghi được xác định bằng hiệu số giữa 2 con trỏ \texttt{\$s4} và \texttt{\$s3}:

\begin{codeblock}[language=MIPS]
  li $v0, 13           # syscall 13: mở file kết quả
  la $a0, file_out     # đường dẫn "formatted_result.txt"
  li $a1, 1            # cờ 1: mở để ghi
  li $a2, 0            # mode bị bỏ qua
  syscall
  bltz $v0, open_error # mở thất bại thì báo lỗi
  move $s0, $v0        # $s0 = descriptor của file kết quả

  li $v0, 15           # syscall 15: ghi file
  move $a0, $s0        # ghi vào file kết quả
  move $a1, $s3        # từ đầu chuỗi kết quả
  sub $a2, $s4, $s3    # số byte = con trỏ ghi - địa chỉ đầu
  syscall

  li $v0, 16           # syscall 16: đóng file kết quả
  move $a0, $s0        # descriptor của file kết quả
  syscall

  j exit
\end{codeblock}

Nếu quá trình mở tệp thất bại, \texttt{syscall 13} trả về giá trị âm trong \texttt{\$v0}; lệnh \texttt{bltz} sẽ chuyển luồng thực thi tới nhãn \texttt{open\_error} để báo lỗi và thoát chương trình:

\begin{codeblock}[language=MIPS]
open_error:        # không mở được file
  li $v0, 4        # syscall 4: in chuỗi
  la $a0, msg_open # in "Error opening file."
  syscall

  j exit
\end{codeblock}

Hai thủ tục phụ trợ \texttt{append\_str} và \texttt{append\_char} cùng sử dụng \texttt{\$s4} làm con trỏ ghi và tự động tăng địa chỉ sau mỗi byte được chép:

\begin{codeblock}[language=MIPS]
append_str:                       # nối cả chuỗi vào đệm kết quả
  lb $t6, 0($a0)                  # đọc ký tự từ chuỗi nguồn
  beq $t6, $zero, append_str_done # gặp ký tự kết thúc thì thôi
  sb $t6, 0($s4)                  # ghi sang chuỗi kết quả
  addi $a0, $a0, 1                # tiến con trỏ nguồn
  addi $s4, $s4, 1                # tiến con trỏ đích
  j append_str
append_str_done:                  # hết chuỗi nguồn
  jr $ra                          # trở về nơi gọi hàm

append_char:                      # nối một ký tự trong $a0
  sb $a0, 0($s4)                  # ghi ký tự vào vị trí con trỏ đích
  addi $s4, $s4, 1                # tiến con trỏ đích
  jr $ra                          # trở về nơi gọi hàm
\end{codeblock}

Kết quả chạy thử nghiệm với tệp \texttt{raw\_input.txt} chứa dữ liệu mẫu của đề bài. Nội dung hiển thị trên màn hình console:

\begin{codeblock}[language={}]
-----Student personal information-----
Name: Davy Jones
ID: 2251234
Address: 168 Ly Thuong Kiet St District 10 Ward 14 HCMC
Age: 69
Religion: None
\end{codeblock}

Tệp \texttt{formatted\_result.txt} được tạo ra với nội dung trùng khớp hoàn toàn:

\begin{codeblock}[language={}]
-----Student personal information-----
Name: Davy Jones
ID: 2251234
Address: 168 Ly Thuong Kiet St District 10 Ward 14 HCMC
Age: 69
Religion: None
\end{codeblock}
